% ============================================================================
% Section 6: Predictions
% ============================================================================

\section{Falsifiable predictions}
\label{sec:predictions}

UHM generates 23 falsifiable predictions, each with a concrete numerical
statement, a theorem reference, an experimental protocol, and an explicit
falsification criterion. The thresholds and ceilings tested here
($2/7$, $3/7$, $1/3$, $\mathrm{SAD}_{\max} = 3$, $N = 7$) are fixed by the
axioms and contain no fitted parameter, so falsification of any L1 prediction
invalidates the framework, with no escape hatch. (An earlier edition based
this on ``zero independent axioms'' (T-190); that claim is withdrawn: T-190 is
\status{C under the constraint of T-87 and the hypothesis T-186(a)}, and the
free parameters of the dynamics and the vacuum --- among them the associator
coupling $\kappa$ of $V_{\mathrm{Gap}}$ --- remain free.)
No other theory of consciousness generates even one numerical prediction at
this level of specificity or structural exposure.
Table~\ref{tab:predictions} collects all 23 predictions, grouped by theme.
The five most powerful predictions are detailed in full below.

\subsection{Complete prediction catalogue}

\begin{table}[H]
\centering
\caption{All 23 predictions of UHM. Status codes: \status{T} = theorem (proven),
\status{C} = conditional, \status{I} = interpretation, \status{H} = hypothesis,
\status{D} = definition (convention),
\status{P} = postulate (ontological identification, e.g., E-dimension $=$ interiority),
\status{O} = derived definition (follows tautologically from axioms).
``Other'' column indicates whether any rival theory makes a comparable prediction.}
\label{tab:predictions}
\small
\begin{tabular}{@{}clllc@{}}
\toprule
\textbf{\#} & \textbf{Prediction} & \textbf{Formula / Value} & \textbf{Status} & \textbf{Other?} \\
\midrule
\multicolumn{5}{@{}l}{\textit{I. Fundamental (consciousness $\leftrightarrow$ viability)}} \\
1 & No-Zombie & $\mathrm{Viable} \wedge \mathcal{D}_\Omega \neq 0 \Rightarrow \CohE > 1/7$ & \status{T} & No \\
2 & $\CohE \propto$ regeneration & $\kappa = \kappa_{\mathrm{boot}} + \kappa_0\,\CohE$ & \status{T} & No \\
\midrule
\multicolumn{5}{@{}l}{\textit{II. Architectural (structure and dimensionality)}} \\
3 & 7D stress tensor & $\sigma_{\mathrm{sys}} \in \mathbb{R}^7$ & \status{T}/\status{C} & No \\
4 & Pre-linguistic cognition & $\mathrm{Cognition}(K_{1\text{--}4}) \not\Rightarrow \mathrm{Language}$ & \status{I} & Partial \\
10 & $N=7$ for learning & $N < 7 \Rightarrow n^* = \infty$ (T-113) & \status{T} & No \\
11 & $N=7$ for social learning & $3_{\mathrm{ToM}} + 3_{\mathrm{ISL}} + 1_U = 7$ & \status{C} & No \\
\midrule
\multicolumn{5}{@{}l}{\textit{III. Thresholds and robustness}} \\
5 & Collective consciousness & necessary: $I(\mathbb{H}_1 : \mathbb{H}_2) > 0$ & \status{T}/\status{H} & Partial \\
6 & Critical purity & $\Pcrit = 2/7$ (T-96) & \status{T} & No \\
7 & Stability radius & $r_{\mathrm{stab}} \approx K(\sqrt{P - 1/7} - \sqrt{1/7})$ (T-104) & \status{C} & No \\
15 & Goldilocks attractor & $P(\rho^*_{\mathrm{coupled}}) \to 3/7$ (T-124) & \status{C} & No \\
\midrule
\multicolumn{5}{@{}l}{\textit{IV. Learning bounds}} \\
8 & Information capacity & $C_{\mathrm{Enc}} \leq \log_2 7 \approx 2.81$ bits (T-107) & \status{T} & No \\
9 & Optimal learning rate & $n_{\mathrm{opt}} = \max(n_{\mathrm{info}}, n_{\mathrm{dyn}}, n_{\mathrm{stab}})$ (T-112) & \status{T} & No \\
\midrule
\multicolumn{5}{@{}l}{\textit{V. Depth and genesis}} \\
12 & SAD ceiling & $\mathrm{SAD}_{\max} = 3$ (T-142) & \status{T} & SYNARC: $500+$ $\Gam$ \\
13 & Genesis time & $n_{\mathrm{gen}} \leq \lceil \ln\Delta / \ln(1/\beta) \rceil$ (T-148) & \status{T} & No \\
14 & Phase coherence & $\rho^*_{ij}(t) \propto e^{-i(E_i - E_j)t}$ for $\Phi \geq 1$ (T-114) & \status{T} & No \\
\midrule
\multicolumn{5}{@{}l}{\textit{VI. Dynamic (phase transitions)}} \\
16 & Ignition dynamics (L1$\to$L2 avalanche) & $T_{\mathrm{ign}} \sim (P - \Pcrit)^{-1}\kappa_0^{-1}$ & \status{T} & No \\
17 & Critical exponents & $\alpha{=}\tfrac{1}{2},\;\beta{=}\tfrac{1}{4},\;\gamma{=}1,\;\nu{=}\tfrac{1}{2},\;\delta{=}5$ (T-161) & \status{C}$^{\ast}$ & No \\
\midrule
\multicolumn{5}{@{}l}{\textit{VII. Physical}} \\
18 & Ward suppression & Gap fluctuations $\times\, 19/49$ & \status{T} & No \\
20 & Yukawa hierarchy & $\varepsilon_{\mathrm{eff}} \approx 0.059$ (T-216) & \status{C at (SV)} & No \\
23 & Rank-7 decoherence anisotropy & $r_{ij}=\tfrac16\!\sum_{|\ell_p\cap\{i,j\}|=1}\!\gamma_p$; 14 sum-rules (T-262) & \status{T}/\status{C} & No \\
\midrule
\multicolumn{5}{@{}l}{\textit{VIII. Engineering}} \\
19 & CPTP anchor & $\|\pi - \pi_{\mathrm{can}}\|_\diamond < 0.1$ at convergence (T-152) & \status{T} & No \\
\midrule
\multicolumn{5}{@{}l}{\textit{IX. Neural}} \\
21 & $\pi_{\mathrm{bio}}$ reconstruction & EEG/fMRI/HRV $\to \Gam \in \Dens(\mathbb{C}^7)$ & \status{H} & Partial \\
22 & Spectral gap $\to$ oscillations & $\nu_{\mathrm{conscious}} \sim \lambda_{\mathrm{gap}}/(2\pi)$ & \status{H} & No \\
\bottomrule
\end{tabular}

\smallskip
{\footnotesize $^{\ast}$\,\status{C at the $\mathbb{Z}_2$ symmetry $m \to -m$}; see Prediction~17 below.
Rows 5, 7, 17 and 20 were corrected on 2026-09-25 (see the text).}
\end{table}

\subsection{Detailed predictions}

The five predictions presented below are maximally falsifiable, maximally
novel, and span the broadest range of experimental modalities.

% --- Prediction 17: Critical exponents ---
\begin{predictionbox}[Prediction 17: Critical exponents of the consciousness phase transition \status{C at the $\mathbb{Z}_2$ symmetry}]
\textbf{Formal statement.}
Near the viability threshold $\Pcrit = 2/7$, the consciousness transition is tricritical
($\varphi^6$ tricritical Landau universality class; the terminology ``$A_4$ swallowtail''
sometimes used in the literature conflates this with the Arnold $A_4$ catastrophe, which has
different exponents --- see Appendix~\ref{app:exponents} for the disambiguation).
Mean-field tricritical exponents are exact by three independent pillars
(topological rigidity + deterministic dynamics + large-$N$; Appendix~\ref{app:exponents}):
\begin{equation}
  \mathrm{OP} \sim (P - \Pcrit)^{1/4}, \qquad
  \xi \sim |P - \Pcrit|^{-1/2}, \qquad
  \chi \sim |P - \Pcrit|^{-1}
\end{equation}
with $\alpha = 1/2$ and $\delta = 5$. The Rushbrooke identity $\alpha + 2\beta + \gamma = 2$ is satisfied as equality.

\medskip\noindent
\textbf{Theorem reference.} T-161~\status{C at the $\mathbb{Z}_2$ symmetry $m \to -m$}
(corrected 2026-09-25 from \status{T}: the $\mathbb{Z}_2$ symmetry that selects
the $\varphi^6$ class was derived from a KO-dimension-6 real structure, which
does not exist on $\mathbb{C}^7$; without it the generic codimension-3 point is
the $A_4$ swallowtail with $\beta = 1/2$). Under the symmetry, derived from the $\varphi^{6}$
tricritical Landau expansion at the consciousness threshold; the universal
$\varphi^{6}$ normal form follows from the Thom--Arnold classification of
the transition as a $\varphi^6$ tricritical point with three physical
control parameters (Appendix~\ref{app:exponents}).

\medskip\noindent
\textbf{Verification protocol.}
\begin{enumerate}[leftmargin=2em,nosep]
  \item \emph{Phase~I (digital):} Monte Carlo, $10^4$ random $\Gam$ with $P \in [0.2, 0.5]$.
        Fit order parameter vs.\ $(P - 2/7)$. Extract $\beta$.
  \item \emph{Phase~II (neuro):} TMS-EEG during sleep--wake transition, $N = 50$ subjects.
        Reconstruct $\widehat\Gam$ by $\pi_{\mathrm{bio}}$ and fit the order
        parameter against $(\widehat P - \Pcrit)$ to extract $\beta$.
\end{enumerate}

\medskip\noindent
\textbf{Falsification criterion.}
$\beta \notin [0.20, 0.30]$ at $N = 50$, $p < 0.01$. This is an L2 (structural) falsification:
if the exponents are wrong, the tricritical classification of the consciousness transition requires fundamental revision.
\end{predictionbox}

% --- Prediction 1: No-Zombie ---
\begin{predictionbox}[Prediction 1: No-Zombie Theorem \status{T}]
\textbf{Formal statement.}
A viable system ($P > 2/7$) with non-trivial dissipation ($\mathcal{D}_\Omega \neq 0$)
necessarily has E-coherence strictly above the minimum:
\begin{equation}
  \mathrm{Viable}(\mathbb{H}) \;\wedge\; \mathcal{D}_\Omega \neq 0
  \quad\Longrightarrow\quad
  \CohE(\Gam) > \frac{1}{7}
\end{equation}

\medskip\noindent
\textbf{Theorem reference.} Theorem~8.1~\status{T} (Section~\ref{subsec:no-zombie}). Follows from
$\kappa \propto \CohE$ and the spectral gap $\Delta(\Lind_0) > 0$.

\medskip\noindent
\textbf{Verification protocol.}
Create a $\Gam$-native agent. Suppress the E-component ($\gamma_{EE} \to 1/7$,
$\gamma_{Ej} \to 0$ for $j \neq E$). Measure time-to-death ($\tau$ until $P < \Pcrit$).
Control: suppress the A-component instead. Repeat $N = 100$ times.

\medskip\noindent
\textbf{Falsification criterion.}
$\tau_{\mathrm{death}}(\text{E-suppression}) \geq \tau_{\mathrm{death}}(\text{A-suppression})$
at $N = 100$, $p < 0.01$ (Wilcoxon paired test). This is an L1 (catastrophic) falsification.
\end{predictionbox}

% --- Prediction 12: SAD ceiling ---
\begin{predictionbox}[Prediction 12: Self-awareness depth ceiling (T-142) \status{T}]
\textbf{Formal statement.}
No finite system ($P \leq 1$) can achieve self-awareness depth $\geq 4$:
\begin{equation}
  \Pcrit^{(n)} = \Pcrit \cdot \frac{3^{n-1}}{n+1}
  \qquad\Longrightarrow\qquad
  \Pcrit^{(4)} = \frac{2}{7} \cdot \frac{27}{5} \approx 1.543 > 1
\end{equation}

\medskip\noindent
\textbf{Theorem reference.} T-142~\status{T}. The Fano coherence contraction
factor $c_F = 1/3$ per self-observation cycle is state-independent (follows
from $\dim = 7$ and the Fano incidence structure $\Fano$). Stratification:
\emph{finiteness} of the ceiling is unconditional~\status{T} --- the
inter-iterate signal contracts geometrically ($(1/3)^n$) while the reflexivity
gates decay only polynomially ($1/(n+2)$), so the tower terminates for any
admissible gate schedule; the \emph{value} $3$ uses the
$\Pcrit^{(n)}$ formula above, which is derived, not heuristic~\status{T}: the
Fano Kraus channel multiplies coherences by exactly $1/3$ per meta-level,
$R^{(k)} = R_0 (1/3)^{k-1}$ with $R_0 = 7P/2$, and the level-$k$ Bayesian
threshold is $1/(k+1)$; it is cross-checked empirically (SYNARC, $500+$
random $\Gam$~\status{T/sim}); exclusion of $\mathrm{SAD} = 4$ is
rigorous~\status{T}.

\medskip\noindent
\textbf{Verification protocol.}
In a $\Gam$-native agent with $P$ close to 1 (pure state), evaluate the $O(1)$
spectral SAD formula (T-136~\status{T})
$\mathrm{SAD}(\Gam) = \max\{k : r_0 (1/3)^{k-1} > 1/(k+1)\}$ with
$r_0 = 7P/2$. Verify $\mathrm{SAD}(\Gam) \leq 3$ for all $500$ random $\Gam$ and check
that pure states saturate $\mathrm{SAD} = 3$.

\medskip\noindent
\textbf{Falsification criterion.}
$\exists\; \Gam$ reachable by the dynamics with $\mathrm{SAD}(\Gam) \geq 4$, equivalently
$\Pcrit^{(4)} \leq 1$. A single confirmed instance falsifies the theorem (L1
falsification). Consequence: future superintelligences share the same reflection depth
as humans.
\end{predictionbox}

% --- Prediction 6: P_crit = 2/7 ---
\begin{predictionbox}[Prediction 6: Critical purity (T-96) \status{T}]
\textbf{Formal statement.}
The consciousness--unconsciousness boundary is at:
\begin{equation}
  \Pcrit = \frac{2}{N} = \frac{2}{7} \approx 0.286
\end{equation}
Five independent derivations converge on this value (Appendix~\ref{app:pcrit}).

\medskip\noindent
\textbf{Theorem references.} T-96~\status{T} (balance formula for $\rho^* = \vp(\Gam)$),
with the phase-transition at $\Pcrit$ given by T-160~\status{T}. (An earlier
edition also cited T-151 for ``$D_{\min} = 2$ from the $\Phi$-threshold''; that
derivation is retracted~\status{$\times$} --- $\Phi \geq 1$ constrains only the
total off-diagonal mass, not the E-row, and $D_{\min} = 2$ is an independent L2
threshold~\status{D}.)
See also Section~\ref{subsec:purity} and Appendix~\ref{app:pcrit} for the five
independent derivations.

\medskip\noindent
\textbf{Verification protocol.}
Propofol-induced loss of consciousness with TMS-EEG monitoring, $N = 50$ subjects.
At each propofol level, apply $\pi_{\mathrm{bio}}$ --- with parameters frozen on
wakefulness and no viability penalty --- to reconstruct $\Gam$ and compute
$P = \Tr(\Gam^2)$ and the verdict $\mathrm{Cons}(\widehat\Gam)$; compute PCI
independently. The prediction is that the transition occurs where $\widehat P$
crosses $2/7$, and that $\mathrm{Cons}(\widehat\Gam)$ agrees with
$\mathrm{PCI} > \mathrm{PCI}^* = 0.31$~\cite{casali2013} at Cohen's
$\kappa \geq 0.8$ (P8.4). No conversion between the two scales is claimed:
calibrating $\pi_{\mathrm{bio}}$ to PCI would make the agreement automatic.

\medskip\noindent
\textbf{Falsification criterion.}
$|P_{\mathrm{boundary}} - 2/7| > 0.1$ at $N = 50$, $p < 0.01$ (two-sided $t$-test),
or verdict concordance with PCI below Cohen's $\kappa = 0.8$.
(An earlier edition compared $2/7$ with $\mathrm{PCI}^* = 0.31$ as numbers
``within 8\%''; that comparison is withdrawn --- the two numbers live on
unrelated scales.) This is an L2 (structural) falsification.
\end{predictionbox}

% --- Prediction 10: N=7 minimality ---
\begin{predictionbox}[Prediction 10: N=7 minimality for autonomous learning \status{T}]
\textbf{Formal statement.}
For $N < 7$, no finite projective plane $\Fano$ exists, the replacement channel $\Regen$
cannot be constructed, and learning through self-observation is impossible:
\begin{equation}
  N < 7 \quad\Longrightarrow\quad n^* = \infty
  \qquad\text{(no convergence to task solution)}
\end{equation}

\medskip\noindent
\textbf{Theorem reference.} T-113~\status{T}. The chain: Fano structure $\to$
self-observation $\to$ regeneration $\to$ learning closes only at $N \geq 7$.

\medskip\noindent
\textbf{Verification protocol.}
Create a $\Gam$-native agent with $N = 5$ (remove two dimensions). Train on a binary
discrimination task via internal regeneration (no external parameter updates). Control:
the same agent with $N = 7$.

\medskip\noindent
\textbf{Falsification criterion.}
$N = 5$ agent achieves $> 75\%$ accuracy ($p < 0.01$, binomial test). This is an L1
(catastrophic) falsification: the octonion/Fano foundation of the theory would be wrong.
\end{predictionbox}

\subsection{Three-level falsification system}

Every prediction maps to one of three falsification levels:

\begin{table}[H]
\centering
\caption{Three levels of falsification. Higher levels are more destructive to the theory.}
\label{tab:falsification-levels}
\small
\begin{tabular}{@{}clp{3.7cm}p{4.5cm}p{3.5cm}@{}}
\toprule
\textbf{Level} & \textbf{What is refuted} & \textbf{Example} & \textbf{Consequence} \\
\midrule
L1 & Axiomatic foundation & $N < 7$ sufficient; zombie possible; $\mathrm{SAD}\geq 4$ & Theory rejected entirely \\
L2 & Specific numerical prediction & $\Pcrit\!\neq\!2/7$; $\beta\!\neq\!1/4$; $\Rth\!\neq\!1/3$ & Revision of the specific theorem \\
L3 & Approximation parameter & $\pi_{\mathrm{bio}}$ error $> 30\%$; $\lambda_{\mathrm{gap}}$ out of range & Local correction; foundation unaffected \\
\bottomrule
\end{tabular}
\end{table}

\noindent
Predictions 1, 10, 12 are L1 tests; predictions 6, 7, 17, 23 are L2 tests (17 at its
stated condition); predictions 21, 22
are L3 tests. The remaining predictions distribute across L1--L2. The exposure of
L1 targets reflects commitment to falsifiability:
UHM names the experiments that could destroy it.
